\documentclass[11pt, a4paper]{article}\usepackage[utf8]{inputenc}\usepackage{amsmath}\usepackage[margin=1in]{geometry}\begin{document}\title{Problem Set 9: Formulasi MILP}\maketitle
\section*{Soal 1 (C3)} Diberikan variabel kontinu $x \ge 0$ dan biner $y$. Formulasikan kendala logika: "Jika $y=0$ maka $x=0$, jika $y=1$ maka $x \le 100$" menggunakan teknik \textit{Big-M}.
\section*{Soal 2 (C4)} Modelkan permasalahan Unit Commitment sederhana untuk 3 generator dengan biaya start-up dan batas daya minimum menggunakan MILP.
\end{document}